% !TEX root = main.tex
% **************************************************
% Settings
% **************************************************
% --------------------------
% Packages
% --------------------------
\usepackage[export]{adjustbox}
\usepackage{blindtext}
\usepackage{booktabs}    % nicer horizontal rules in tables
\usepackage[small,hang,bf]{caption}    % caption options
\usepackage{comment}
\usepackage{csquotes}
\usepackage{emptypage}  % no headers in empty pages
\usepackage{eso-pic}    % Para tener control sobre los márgenes y fondo
\usepackage{enumitem}
\usepackage{titlesec}   % Control sobre la apariencia de "Capítulo, Sección, etc"
\usepackage{fancyhdr}   % personalizar encabezado
\usepackage{float}      % control más fino sobre dónde se colocan las imágenes/tablas
\usepackage[T1]{fontenc}
\usepackage[a4paper]{geometry}
\usepackage[acronym,nonumberlist,nomain]{glossaries}  % package for list of acronyms
\usepackage{graphicx,graphics}   % packages to manage images
\usepackage[protrusion=true, expansion=true, nopatch=footnote]{microtype}  % Use microtype to improve readability
\usepackage{lmodern}     %La fuente actual
\usepackage{longtable}   % tables spanning more than one page
\usepackage{ltablex}
\usepackage{listings}
\usepackage{makecell}   % Super útil para arreglar problemas en las tablas, recomiendo usar -P
\usepackage{multicol}    % multi-columns in tables
\usepackage{multirow}    % multi-rows in tables
\usepackage{nicematrix}
\usepackage{setspace}
\usepackage{subcaption}
\usepackage{parskip}     % paragraph spacing (suppress indentation)
\usepackage{ragged2e}
%\usepackage{pdfpages}
\usepackage{tabularx}    % tables spanning the full text width
\usepackage{tikz}                % drawing custom shapes
\usepackage{textcomp}   % Inserting special text characters
\usepackage[table, dvipsnames]{xcolor}
%%% BIBLIOGRAPHY
% The style used is IEEE. A list of available styles is available  at https://www.overleaf.com/learn/latex/Biblatex_bibliography_styles
\usepackage[square,numbers]{natbib}

\usepackage[spanish,english,es-tabla]{babel}
\usepackage{amsmath,amssymb}
\usepackage{hyperref}

% ==================================================
% Customize document
% ==================================================
% Define the geometry of the document
\newgeometry{top=2.5cm, bottom=2.5cm, right=2.5cm, left=2.5cm}

\titleformat{\chapter}[hang]
{\normalfont\LARGE\bfseries} % 1. Formato general (fuente, tamaño, negrita)
{\thechapter.}              % 2. Etiqueta (ej: "1.")
{20pt}                      % 3. Separación horizontal entre el número y el título
{\Huge}                     % 4. Formato específico antes de imprimir el título

% Redefinir los espacios en blanco
\titlespacing*{\chapter}{0pt}{-30pt}{20pt}
\setlength{\parskip}{7pt}   % Change the default value of paragraph spacing
\setlength{\parindent}{0pt}    %Change default paragraph indentation

% ----------------------------
% hiperlinks
% ----------------------------
\hypersetup{
	colorlinks=true,          % Enable colored links
	linkcolor=black,          % Color for internal links
	urlcolor=blue,            % Color para las URL's u texto que mande a URL's
	citecolor=blue            % Color para las citas (comando \cite{} o \citet{})
}

% ----------------------------
% Header and footer
% ----------------------------
\fancyhf{}   % Clear the default header and footer
\renewcommand{\headrulewidth}{0.2pt}

\setlength{\headheight}{14.5pt}
\fancyhead[RO]{\nouppercase \leftmark}  % Define the header for rigth (odd) pages
\fancyhead[LE]{\Asignatura}           % Define the header for left (even) pages
\fancyfoot[C]{\thepage}                 % Define the page number as the footer

% ----------------------------
% Colores para el documento
% ----------------------------
\definecolor{coverBlue}{HTML}{0070C0}
\definecolor{tableHeader}{rgb}{0.7,0.7,0.7}

% ----------------------------
% Español - Las traducciones se pueden cambiar a gusto del autor
% ----------------------------
% Cambia el título del índice general:
\addto\captionsspanish{\renewcommand{\contentsname}{Índice General}}
% Cambia el título de la Bibliografía:
\addto\captionsspanish{\renewcommand{\bibname}{Referencias}} % Alternativamente: Bibliografía
% Cambia los apéndices por anexos:
\addto\captionsspanish{\renewcommand{\appendixname}{Apéndice}} %Deshabilitado porque si son de elaboración propia se deben llamar apéndices
\addto\extrasspanish{\def\appendixautorefname{\appendixname}}
% Cambia capítulo por parte:
\addto\captionsspanish{\renewcommand{\chaptername}{Parte}}

% ----------------------------
% Increase indentation in lists - No he sido capaz de que esta parte del código funcione correctamente, lo dejo a modo de reliquia - P
% ----------------------------
\newlist{mydescription}{description}{1}
\setlist[mydescription,1]{labelindent=2em, leftmargin=2em}
\setlist[description]{align=right, labelwidth=2cm}

% ----------------------------
% Increase height of table cells - 1 no hace nada, si es necesario aumentarlo, cambiarlo por 1.1 e ir incrementando poco a poco. Es un factor de escala, así que cuidado.
% ----------------------------
\renewcommand{\arraystretch}{1}

%------------ Fuentes en las figuras --------------
\newlength{\captionlabelwidth}

\newcommand{\source}[1]{
	\settowidth{\captionlabelwidth}{\small\bfseries\figurename~\thefigure:\hspace{\labelsep}}
	\par\nopagebreak\vspace{2pt}
	\noindent
	\hspace*{\captionlabelwidth}
	\parbox[t]{\dimexpr\linewidth-\captionlabelwidth\relax}{
		\small FUENTE: #1
	}
}
% --------------- TODO Command Configuration ---------------
\makeatletter % Allow use of @ in command names

\newcounter{mytodocounter}

\newcommand{\todo}[1]{%
	\stepcounter{mytodocounter}%
	\ifnum\value{mytodocounter}=1 % Only do this for the first todo item
		\@bsphack % Prevents unwanted spaces
		\protected@write\@auxout{}{%
			\string\gdef\string\todosfound{true}% Define \todosfound globally for next run
		}\@esphack % Prevents unwanted spaces
	\fi
	\textcolor{red}{\textbf{ToDo~\themytodocounter:} #1}%
	\addcontentsline{tdo}{todoitem}{\protect\numberline{\themytodocounter}#1}%
}

\newcommand{\listoftodos}{%
	\@ifundefined{todosfound}{}{
		\chapter*{Lista de cosas pendientes}% 
		\addcontentsline{toc}{chapter}{Lista de cosas pendientes}
		\@starttoc{tdo}%
	}%
}

\newcommand*\l@todoitem[2]{%
	\setlength{\parskip}{7pt}
	\par
	\setlength{\@tempdima}{1.5em}% 
	\noindent
	{#1}% 
	\nobreak
	\hfill
	(#2)%
	\par
}

\makeatother % Revert @ to its special meaning
% --------------- End of TODO Command Configuration ---------------
